\documentclass[master=ecws,masteroption=ss]{kulemt}
\setup{
  title={Constraint Generation with Parametric Higher-Order Abstract Syntax},
  author={Jens Nys},
  promotor={Prof.\,dr.\,ir.\ Devriese},
  assessor={},
  assistant={Dr. Ir.\ S. Keuchel \and Ir. D. Carnier}}
% Remove the "%" on the next line for generating the cover page
%\setup{coverpageonly}
% Remove the "%" before the next "\setup" to generate only the first pages
% (e.g., if you are a Word user).
%\setup{frontpagesonly}

% If you want to include other LaTeX packages, do it here. 

% Finally the hyperref package is used for pdf files.
% This can be commented out for printed versions.
\usepackage[pdfusetitle,colorlinks,plainpages=false]{hyperref}
\usepackage[table]{xcolor} % 'table'optie is handig voor cell backgrounds
\usepackage{xltabular}    % combinatie van longtable + tabularx
\usepackage[table]{xcolor}
\usepackage{enumitem}
\usepackage{xltabular,enumitem} 
\usepackage[table]{xcolor} 
\usepackage{pdflscape}

\usepackage{minted} % code highlighting
%%%%%%%
% The lipsum package is used to generate random text.
% You never need this in a real master's thesis text!
\IfFileExists{lipsum.sty}%
 {\usepackage{lipsum}\SetLipsumDefault{11-13}}%
 {\newcommand{\lipsum}[1][11-13]{\par And some text: lipsum ##1.\par}}
%%%%%%%
\usepackage{longtable}
%%%%%%%
%\includeonly{chapter-n}
\begin{document}

\begin{preface}
This is my acknowledgment to thank everyone who has kept me busy.
I would like to thank my supervisor, my advisor, and the entire jury.
My family has also been very supportive, of course.
\end{preface}

\tableofcontents*

\begin{abstract}
  This \texttt{abstract} environment provides a summary of the work, which may or may not be extensive.
The intention is that this should be limited to 1 page.

  \lipsum[1]
\end{abstract}

% A list of figures and tables is optional
%\listoffigures
%\listoftables
% If there are only a few figures and tables, it is better to use the following:
\listoffiguresandtables
% The list of symbols is also optional.
% This list must be created manually, e.g. as follows:
\chapter{List of abbreviations and symbols}
\section*{abbreviations}
\begin{flushleft}
  \renewcommand{\arraystretch}{1.1}
  \begin{tabularx}{\textwidth}{@{}p{12mm}X@{}}
    LoG   & Laplacian-of-Gaussian \\
    MSE   & Mean Square error \\
    PSNR  & Peak Signal-to-Noise ratio \\
  \end{tabularx}
\end{flushleft}
\section*{Symbols}
\begin{flushleft}
  \renewcommand{\arraystretch}{1.1}
  \begin{tabularx}{\textwidth}{@{}p{12mm}X@{}}
    42    & ``The Answer to the Ultimate Question of Life, the Universe,
            and Everything'' volgens de \cite{h2g2} \\
    $c$   & Speed of light \\
    $E$   & Energy \\
    $m$   & Mass \\
    $\pi$ & The number pi \\
  \end{tabularx}
\end{flushleft}

% Now comes the main text
\mainmatter

\chapter{Introduction}
\label{cha:intro}
General introduction of the work. Explain on high level the idea of the thesis and the main contributions.


\chapter{Background}
\label{cha:1}
I will intruduce concepts like foas, hoas, phoas and verification condition generation. Then I will motivate using phoas for verification condition generation

\section{Representing abstract syntax}

\subsection{First Order Abstract Syntax}
Here I will describe the first order approach for representing abstract syntax, I can describe nominal string representations/de bruijn levels and mention their strenghts and weaknesses. (information from eg. https://jesper.sikanda.be/posts/1001-syntax-representations.html)
\subsection{Higher order abstract syntax}
Describe HOAS representation 



\subsection{Parametric Higher order abstract syntax}

(information from eg. Parametric Higher-Order Abstract Syntax for Mechanized
Semantics)

\subsection{Translations between representations}
Describe translations phoas<->foas and phoas-> hoas such as described in Unembedding Domain-Specific Languages.


\section{monads?}
(I'm not sure if this is necessary but it is something most other students would not have as background knowledge)
Describe option, reader, state and continuation monads.

\section{Verification condition generation}
Describe what a verification condition is and how it can be computed by recursively calculating the weakest precondition.

\subsection{Wstore and its combinators}
Describe the Wstore monad and assert, demonic and angelic choice combinators.
\subsection{generating verification conditions}
Describe how to use these combinators to generate the weakest precondition.



(Verified Symbolic Execution with Kripke Specification
Monads (and No Meta-programming)



\chapter{Verification condition generation}
\label{cha:2}
\textit{I will introduce the syntax and semantics of the programming language that I am writing the verification condition generation for. I will then describe the first order and higher order proposition languages and contracts written in them. Finally I will explain how the verification condition generation works.} 

When writing a verification constraint generator you need some ingredients as input. We want to generate a formula that asserts whether a program satisfies its specification: executing the program from a state in which the precondition holds, always ends in a state in which the postcondition holds. This means we have to have a language in which the program is written, but we also need a language in which the pre and postconditions are written. For both of these, we both need a syntax and a semantics.


\section{Programming language}
\textit{Describe the syntax and semantics of the programming language of my vfc generator.}



\subsection{Syntax}
To describe the syntax of the programming language, we describe the datatype we use to represent it. A Program consists of a function that can take various input variables and then executes a statement that mentions these variables. Statements are either expressions, variable introduction with Let, variable reassignment with Assign, the executions of 2 statements sequentially after one another or a conditional control flow branch. Our language only has Integer values, so expressions are either integer literals, variables or integer operations of 2 expressions.


\begin{minted}{haskell}

--program variables
type X = String
--function names
type F = String

type Value = Int

data Exp = Lit Value
    | Var X
    | Add Exp Exp
    | Mul Exp Exp
    | Minus Exp Exp
  deriving (Eq,Show)
  
data Stm = Expr Exp
    | Assign X Exp -- x := Stm (update a variable)
    | Let X Exp Stm -- let X = Stm where Stm (make a new variable)
    | Seq Stm Stm --e1;e2
    | If Bexp Stm Stm -- if bexp then stm else stm
  deriving (Eq,Show)

data Relop = Equal | LessThan | GreaterThan | LessThanEqual | GreaterThanEqual

data Prog = Fun F [X] Stm
\end{minted}


\subsection{Semantics}
To describe a meaning to this syntax, we can give an interpreter for this language.

\begin{itemize}
    \item State: State: Consider the statement (Var “x”). It should have a different result if a previous assignment is made or not. We need an environment that keeps track of this state. This could lead to a function of type Stm->Env->Int. Since we sometimes can edit the environment, such as in an Assign statement, we could also return a tuple with the updated environment Stm->Env->(Int,Env). This can be made easier to work with by using a state monad. Then the return type can be Stm -> State Env Int
    \item Failure (Maybe): If we come across a statement (Var “x”) but there is no “x” in the environment, what should we return? We add the possibility of failure to the result by returning a Maybe Int. Combining with the state monad by using a stateT transformer, the type for the interpreter becomes  Stm -> StateT (Store Value) Maybe Value
\end{itemize}

\section{Proposition languages}
\textit{Describe the syntax of the first order, higher order and parametric higher order versions of the proposition languages. }
\subsection{syntax}
The proposition language of our language should represent first order logic expressions in which we have a way of comparing expressions.

\subsubsection{HOAS propositions}
A HOAS representation for these propositions would look like the following:

\begin{minted}{haskell}
data Prop = T 
    | F 
    | Not Prop
    | And Prop Prop 
    | Or Prop Prop 
    | Implies Prop Prop
    | Cmp Relop Value Value 
    | Exist (Value->Prop) 
    | Forall (Value->Prop)
\end{minted}
The unary and binary combinators are straight forwared. Variable introductions in Exist and Forall use the HOAS approach. The place where these introduced variables can occur are in the comparisons in the comparisons. Notice that our program still typechecks if we write a formula like this
\begin{minted}{haskell}
Forall (\x -> Cmp SmallerThan x (x+1))  
\end{minted}
We can write meta language "expressions" in our comparisons. This property is something we will want in our first order and parametric higher order propositions as well. 

\subsubsection{phoas propositions}
The phoas representation for propositions I'll be using has te following syntax:
\begin{minted}{haskell}
data Prop v = T 
    | F 
    | Not Prop
    | And Prop Prop 
    | Or Prop Prop 
    | Implies Prop Prop
    | Cmp Relop v v 
    | Exist (Maybe String) (v->(Prop v))
    | Forall (Maybe String) (v->(Prop v))
\end{minted}

Notice that if we ignore the Maybe String parameter for the quantifiers and we were to instantiate v with Value we would get a type equivalent to the HOAS propositions. We still have the advantage of having well-formedness . There are however some problems with the approach. 


For example, if we don't assume anything else about the type v above, our comparisons could only compare variables against one another. We would not be able to compare if a variable is greater than 0. Therefore, in most of what is to come, we will expect that the type v has implemented the following typeclass:

\begin{minted}{haskell}
class ValueAlgebra v where
    lit :: Value -> v
    add :: v -> v -> v
    minus :: v -> v -> v
    mul :: v -> v -> v
\end{minted}

This way, if we assume the typeclass ValueAlgebra v we can write propositions like 
\begin{minted}{haskell}
Forall (\x -> Cmp SmallerThan x (add x (lit 1)))  
\end{minted}

Here we can compare expressions of variables instead of only variables.


Another problem with this phoas approach is that printing out formulas or reasoning about equality of formulas is not easy. Also in the other direction, this way of writing propositions restricts us to using haskell to state the propositions since we use haskell's variable binding construct. Therefore it would be nice to have a first order representation that we could use to state our contract and read printed propositions from. Additionally it would be nice to be able to translate to and from these first order representation. Luckily, this is possible, as I will demonstrate in \ref{translation}. Since these translations will have to pick a string to represent the variables that are represented with meta-language variables here, we might want to "suggest" what the name of a variable might be if it were to be translated. Therefore, the quantifiers of our phoas type have an additional field in their constructor which we can be used to suggest a name for the introduced variable in case we were to translate the proposition. 

\section{first order propositions}
The first order representation for propositions I'll be using has the following syntax:

\begin{minted}{haskell}
data Prop v = T 
    | F 
    | Not Prop
    | And Prop Prop 
    | Or Prop Prop 
    | Implies Prop Prop
    | Cmp Relop Stm Stm 
    | Exist string FoasProp 
    | Forall string FoasProp
\end{minted}

Variables are represented as strings and comparisons are represented as expressions with 

Notice that this approach allows for ill-formed propositions. For example, we could construct the following ill-formed proposition where we use "y" as a variable that is not introduced:
\begin{minted}{haskell}
Forall "x" (Cmp Equal "x" "y")
\end{minted}




\subsection{Translations \label{translation}}
\textit{Describe the translations between Foas and PHOAS propositions}
\subsubsection{From phoas to foas propositions}

To translate a phoas proposition to a foas proposition, we can pick an instanitation for our type parameter that can help us in doing the translation. Since the foas comparisons happen between expressions and the phoas representation compares values of our type parameter, we can instantiate our type parameter with Exp. This makes our comparison case, where we use variables straight forward, since the expressions to compare will already be there.

The introduction of variables in Forall and Exist is a little more involved. We need a way to come up with fresh variable names for the phoas variables that don't have  variable identity. These names should be fresh because otherwise variable shadowing might change the semantic meaning of a proposition. For example, if we were to pick every variable name the same as "x", then the formula "Forall Nothing (a -> Forall Nothing (b -> Cmp Equal a b)) (which is not true for values a:=1, b:=2) would translate to the formula "Forall "x" (Forall "x" (Cmp Equal (Var "x") (Var "x"))" which is a tautology. 

To pick these names fresh we pass a integer parameter that we increment at each quantification so we can name our variables "x0", "x1", "x2",… if no name suggestion is given and we pick the suggestion if one is given. To avoid undefined behaviour, we assume that the suggestions don’t have the form “x Integer”. For this extra integer parameter, we can use a Reader monad to write our implementation using do notation. This gives us the following type signature.


\begin{minted}{haskell}
phoas_to_foas_reader :: PhoasProp Exp -> ReaderInt FoasProp
\end{minted}


We can obtain the translation from this by letting our integer parameter start with 0.

\begin{minted}{haskell}
phoas_to_foas :: PhoasProp Exp ->  FoasProp
phoas_to_foas phoasProp = runReader (phoas_to_foas_reader phoasProp) 0
\end{minted}

Once we have a name for our variable, for example "x0", we can construct the expression (Var "x0") and provide it to our higher order function. This gives us back a phoas proposition, which we can recursively translate to a foas proposition. The code for this is the following:

\begin{minted}{haskell}
    do i <- ask
      let arg = "x" ++ show i
    
      case m of
        Nothing -> do --no suggestion provided, we pick arg as our fresh variable and increment i with 1.
           body <- (local (+1) $ phoas_to_foas_reader $ (f (Var arg)))
           return $ quantifier arg body
    
    
    
        Just s -> do -- a suggestion was provided, we use the suggestion.
           body <-              (phoas_to_foas_reader $ (f  (Var s)))
           return $ quantifier s body
\end{minted}

\textit{Perhaps the following simplified version is more readable by skipping the suggestion.}
\begin{minted}{haskell}
    do i <- ask
      let arg = "x" ++ show i
    
       do 
           body <- (local (+1) $ phoas_to_foas_reader $ (f (Var arg)))
           return $ Forall arg body
\end{minted}

By using the reader monad, we obtain locally different names within a scope. For example, consider the following phoas proposition:

\begin{minted}{haskell}
PhoasAnd (PhoasForall Nothing (\a -> PhoasCmp Equal a a)) (PhoasExist Nothing (\a-> PhoasCmp Equal a a ))
\end{minted}

This will be translated to the following foas proposition:


\begin{minted}{haskell}
(And (Forall "x0" (Cmp Equal (Var "x0") (Var "x0"))) (Exist "x0" (Cmp Equal (Var "x0") (Var "x0"))))
\end{minted}

While these introduced variables have the same name, there will never be 2 variables with the same name in the same scope. Shadowing is therefore not a problem.


However, if we want all introduced variables to have a different name, even if they are in different scopes we could use a state monad instead of the reader monad. In that case the example would be translated to
\begin{minted}{haskell}
(And (Forall "x0" (Cmp Equal (Var "x0") (Var "x0"))) (Exist "x1" (Cmp Equal (Var "x1") (Var "x1")))).
\end{minted}


\subsubsection{From foas to phoas propositions}

The translation from a foas to a phoas proposition comes with some hickups. Since foas propositions can be ill-formed and phoas propositions can't, there are some foas propositions for which the translation will fail. This failure happens when a variable is used that is not introduced. The translation makes use of a helper function that maps introduced variable names to the logical variable we use to represent it in the phoas proposition. the translation from foas to phoas starts by having this map as empty and uses a helper function with the following type

\begin{minted}{haskell}
foas_to_phoas' :: ValueAlgebra a => FoasProp -> Map String a -> PhoasProp a


\end{minted}


The interesting cases happen when a variable is introduced and when it is used, as seen below. Notice how we introduce a meta language variable in the Forall case and put it in the store. When a variable is used in forall, we have to lift the expression of type Exp into an element of the arbitrary type. To do this, we interpret the expression into the arbitrary type by requiring the valueAlgebra interface to create expressions within that type.


\begin{minted}{haskell}
    
foas_to_phoas' (FoasForall lvar p) env = PhoasForall (Just lvar) (\v -> foas_to_phoas' p (insert lvar v env))
foas_to_phoas' (FoasCmp relop s1 s2) env = PhoasCmp relop (expression_to_algebra s1 env) (expression_to_algebra s2 env)
\end{minted}











\section{Generating verification conditions}






Describe where my implementation differs from the hoas constraint generation described in the background chapter.

If we want to verify a hoare triple {Pre} program {Post}, the postcondition usually says something about the obtained result. Lets say a program returns a type r, then our postcondition would have the form (r->Prop) where Prop is a proposition. The precondition can be any Proposition. Our constraint generator will work as follows: We get our program and a postcondition for our program. We will then, for every instruction, calculate the weakest precondition wp(program,post). This is the precondition such that for every other precondition Pre for which the hoare triple {Pre} program {Post} holds, Pre implies wp(program,post). Our verification condition will then be: “for all possible inputs, does our precondition imply wp(program(Input),post)”


\subsection{}


\chapter{Formal properties}
\label{cha:3}
I will describe and a provide a proof that the provided approach actually leads to a well-scoped verification condition.I will also provide a proof that the verification condition is adequate against the semantics of the language.

\section{Mechanization}
I will first describe some of the changes to support . These include simpler language, big step semantics and a "default" in the foas->phoas translation (haskell code errors if the foas proposition is ill-formed, coq picks "lit 0").

\section{well-formedness proof}
Describe end to end theorem of well-formedness and the major lemma's involved



\section{Conclusion}


\chapter{Evaluation}
\label{cha:4}


\section{Quantitative}
Benchmark haskell constraint generation, maybe it's interesting to benchmark different instantiations of type parameter in the translation to foas

Maybe write a first order constraint solver and compare against that?

compare lines of code.

\section{Qualitative}

Describe advantages of approach
* able to use library methods because we don't work in "indexed category"

No variable bookkeeping for programming

In possible future proof assistants might have parametricity results available for free, which would reduce the proof effort of the well-formedness proof.

In adequacy proof, no "code duplication" for the relation between foas and phoas instantiation.





% ... and so on until
%\chapter{The Final Chapter}
\label{cha:n}
\lipsum[79]

\section{The First Topic of this Chapter}
\subsection{Item 1}
\subsubsection{Sub-item 1}
\lipsum[80]

\subsubsection{Sub-item 2}
\lipsum[81]

\subsection{Item 2}
\lipsum[82]

\section{The Second Topic}
\lipsum[83-85]

\section{Conclusion}
\lipsum[86-88]

\chapter{Conclusion}
\label{cha:conclusion}
The final chapter contains the overall conclusion. It also contains
suggestions for future work and industrial applications.

\lipsum[1-7]


% If you have appendices:
\appendixpage*          % if wanted
\appendix
\chapter{The First Appendix}
\label{app:A}
Appendices hold useful data which is not essential to understand the work
done in the master's thesis. An example is a (program) source.
An appendix can also have sections as well as figures and references\cite{h2g2}.

\section{More Lorem}
\lipsum[50]

\subsection{Lorem 15--17}
\lipsum[15-17]

\subsection{Lorem 18--19}
\lipsum[18-19]

\section{Lorem 51}
\lipsum[51]

% ... and so on until
\chapter{The Last Appendix}
\label{app:n}
Appendices are numbered with letters, but the sections and subsections use
arabic numerals, as can be seen below.

\section{Lorem 20-24}
\lipsum[20-24]

\section{Lorem 25-27}
\lipsum[25-27]

\backmatter
% The bibliography comes after the appendices.
% You can replace the standard "abbrv" bibliography style by another one.
\bibliographystyle{abbrv}
\bibliography{references}
\chapter*{Code of conduct and transparency statement on the use of GenAI for KU Leuven-students (academic year 2025--2026)}
\addcontentsline{toc}{chapter}{Code of Conduct}
\markboth{Code of Conduct}{Code of Conduct}
Generative AI (GenAI) assistance tools can be used to generate various types of content, including text, images, code, video, music, or combinations thereof. Common examples of such tools include ChatGPT, Google Gemini, Microsoft Copilot, Midjourney, Claude.ai, Perplexity.ai, and DALL·E, among others.
\newline
\newline
This code of conduct is a tool that helps students to be transparent about the use of GenAI and fits within the university’s principles on academic integrity.

\subsection*{Important guidelines and remarks}

\begin{itemize}
    \item \textbf{Sensitive or personal data:} Some GenAI tools protect your input and sensitive or personal data better than others. There is often no transparency on what the owners of the AI applications do with the data entered. Therefore, \textbf{do not enter sensitive or personal data in free GenAI tools.} More info about what sensitive and personal data are, is described in the \href{https://www.kuleuven.be/english/education/leuvenlearninglab/support/toolguide/guidelines-for-safe-use-of-genai-tools}{three confidentiality levels of the KU Leuven data classification model} (non-confidential, confidential, strictly confidential). For confidential data you should preferably use \textbf{M365 Copilot Chat with your KU Leuven account}. If you use another GenAI tool, you must be \textbf{absolutely certain} it does \textbf{not store or reuse the data you enter. Try to avoid entering these data. Do so only if strictly necessary, and only to the extent required. For personal data, work with anonymous or pseudonymized data.} Additionally, in case of strictly confidential data this should first be discussed with the teaching staff of the course or your thesis supervisor.
    
    \item \textbf{Copyrighted materials:} For \textbf{lawfully obtained copyrighted material} you should preferably use M\textbf{365 Copilot Chat with your KU Leuven account.} If you use another GenAI tool, you must be \textbf{absolutely certain} it does n\textbf{ot store or reuse the data you enter.}
    
    \item GenAI assistance may not be used for data or topics covered by a \href{https://www.kuleuven.be/english/stuvo/student-legal-services/research-with-interaction-with-a-company-or-organisation}{Non-Disclosure Agreement (NDA)}. Even in the absence of an NDA, certain information may still need to be treated as confidential—for example, due to regulatory requirements or the risk of significant harm to the university if disclosed. In case of doubt, check with your teaching staff or supervisor.
    
    \item If your master’s thesis is under \href{https://iiw.kuleuven.be/english/students/master-thesis/embargo}{embargo}, you should first discuss with your supervisor whether the use of GenAI is permitted.
    
    \item Before using a GenAI tool, always consider whether its use is responsible, including from a \href{https://www.kuleuven.be/english/genai/sustainable-and-responsible-use-of-genai}{sustainability perspective} (e.g. when using GenAI as a search engine, language assistant,…).
    
    \item \textbf{Take a scientific and critical attitude} when interacting with GenAI assistance and interpreting its output, that may not always be correct. 
    
    \item As a student you are responsible for complying with Article 84 of the Regulations on Education and Examinations: your report or thesis should reflect your own knowledge, understanding and skills. Be aware that plagiarism rules also apply to (work that is the result of) the use of GenAI assistance tools. 
    \newline
    \textbf{Exam Regulations Article 84:} \textit{“Every conduct individual students display with which they (partially) inhibit or attempt to inhibit a correct judgement of their own knowledge, understanding and/or skills or those of other students, is considered an irregularity which may result in a suitable penalty. A special type of irregularity is plagiarism, i.e. copying the work (ideas, texts, structures, designs, images, plans, codes , ...) of others or prior personal work in an exact or slightly modified way without adequately acknowledging the sources. Every possession of prohibited resources during an examination (see article 65) is considered an irregularity.”}
    
    \item In order to maintain academic integrity and avoid plagiarism, \textbf{more information about being transparent on the use of GenAI assistance and about correctly citing and referencing GenAI} can be found on this website for students (\href{https://www.kuleuven.be/onderwijs/student/onderwijstools/artificiele-intelligentie#Transparantie}{Dutch}/\href{https://www.kuleuven.be/english/education/student/educational-tools/generative-artificial-intelligence#Transparency}{English}). 
    
    \item \textbf{Additional reading: KU Leuven guidelines on responsible use of Generative AI tools, and other information }(\href{https://www.kuleuven.be/genai}{Dutch}/\href{https://www.kuleuven.be/english/genai}{English})
\end{itemize}

\subsection*{A few final words}

\textbf{If you are uncertain whether or not you should declare your use of GenAI tools, we suggest that you discuss this with your instructor or supervisor. It is always safer to declare GenAI use, even when it is not strictly required. However, declaring GenAI use does not entail that its use is allowed; the right column in the table below provides more detailed instructions in this regard (code of conduct). }
\newline
\newline
Moreover, advanced AI tools are evolving rapidly, and their capabilities have expanded significantly in a short period of time. As a result, we do not yet have all the answers about their responsible use. Finally, it is important to follow-up on the most recent evolutions in AI technologies, to have an open mind but also to be a bit cautious, to communicate with instructors, teaching assistants, supervisors and peers, to be as transparent as we can, and to learn together as we move along.
\newpage
% --- Student info row ---
% --TO BE FILLED IN--
\noindent
\textbf{Student name:} \dotfill \hfill \textbf{Student number:} \dotfill

\vspace{1em}

% --- Grey box section title ---
\noindent
\colorbox{lightgray}{\parbox{\dimexpr\linewidth-2\fboxsep\relax}{\textbf{Please indicate with "X" whether it relates to a course assignment or to the Bachelor’s or Master’s thesis:}}}

\vspace{1em}

O This form is related to a \textbf{course assignment.}

\begin{tabularx}{\linewidth}{@{}lXlX@{}}
\textbf{Course name:} & \dotfill & \textbf{Course number:} & \dotfill \\
\end{tabularx}

\vspace{1em}

This form is related to my O \textbf{Bachelor’s} or O \textbf{Master’s thesis.}

\begin{tabularx}{\linewidth}{@{}lXlX@{}}
\textbf{Title Bachelor’s or Master’s thesis:} & \dotfill & \textbf{Supervisor:} & \dotfill \\
\end{tabularx}

\vspace{1em}

\textbf{Daily supervisor:} \dotfill

\vspace{2em}

% --- Grey box section title ---
\noindent
\colorbox{lightgray}{\parbox{\dimexpr\linewidth-2\fboxsep\relax}{\textbf{Please indicate with "X":}}}

\vspace{1em}

O I \textbf{did not use} any GenAI assistance tool. \\[1em]

O I \textbf{did use} GenAI Assistance. In this case \textbf{specify which ones} (e.g.\ ChatGPT, M365 Copilot, …): \dotfill

\newpage
% ====== PREAMBLE AANVULLINGEN ======

\definecolor{lightgray}{gray}{0.9}

% Kolomdefinitie: mooie linkslijnende X-kolom
\newcolumntype{Y}{>{\raggedright\arraybackslash}X}

% Veilige cel-wrapper (geen \setlist hier!)
\newcommand{\Cell}[1]{\begin{minipage}[t]{\linewidth}\raggedright\small #1\end{minipage}}

% (optioneel) iets compactere tabel
\setlength{\tabcolsep}{6pt}
\renewcommand{\arraystretch}{1.2}

\begin{longtable}{|p{0.25\textwidth}|p{0.2\textwidth}|p{0.55\textwidth}|}
\hline 
\textbf{GenAI assistance used as/for:} & \textbf{Name of the GenAI tool(s) used} If helpful, also describe in which way you were using GenAI related to what is specified as code of conduct. & \textbf{Code of conduct: \newline
For each of the categories below, always take into account the important guidelines and remarks mentioned above (e.g. copyrighted data, sensitive or personal data,…).} \\ \hline
\endfirsthead

%\hline
%\textbf{GenAI assistance used as/for:} & \textbf{Name of the GenAI tool(s) used} If helpful, also describe in which way you were using GenAI related to what is specified as code of conduct. & \textbf{Code of conduct:} For each of the categories below, always take into account the important guidelines and remarks mentioned above (e.g. copyrighted data, sensitive or personal data,…). \\  \\ \hline
%\endhead

\hline
\endfoot

\hline
\endlastfoot

\textbf{As a language assistant for reviewing or improving texts I wrote myself} & \textit{add text here} & This use is similar to using spelling and grammar check tools. In general, you do not have to refer to such kind of GenAI use in the text. 

\textbf{However, be careful: }
\begin{itemize}
    \item When using GenAI tools on texts you did not write yourself to improve the text, \textbf{you have to \href{https://bib.kuleuven.be/english/training-and-tutorials/citation}{refer} to the original source or author}, otherwise you are committing \href{https://www.kuleuven.be/english/education/plagiarism/index}{plagiarism} and thus an irregularity.
\end{itemize} \\ \hline
\textbf{As a paraphrasing tool} & \textit{add text here} & This use is allowed except when it is prohibited by the teacher or the program of study. You may paraphrase your own text or texts by an author other than yourself and take inspiration from what a GenAI tool or another tool suggests (unless it is not allowed). In general, you do not have to refer to such kind of GenAI use or other paraphrasing tools in the text. 

\textbf{However, be careful}: 
\begin{itemize}
    \item If it entails text by an author other than yourself, \textbf{you are not allowed to include that paraphrased text without \href{https://bib.kuleuven.be/english/training-and-tutorials/citation}{reference} to the original source or author}. Without such reference, you would be committing \href{https://www.kuleuven.be/english/education/plagiarism/index}{plagiarism} and thus an irregularity.
\end{itemize} \\ \hline
\textbf{For translation aid to improve texts I wrote myself or to better understand text from others} & \textit{add text here} & This use is allowed except when it is prohibited by the teacher or the program of study. It is similar to using translation tools (Google translate, DeepL, …). In general, you do not have to refer to such kind of GenAI use in the text. 

\textbf{However, be careful:} 
\begin{itemize}
    \item \textbf{You are not allowed to include that translated text without \href{https://bib.kuleuven.be/english/training-and-tutorials/citation}{reference} to the original source or author.} Without such reference, you would be committing \href{https://www.kuleuven.be/english/education/plagiarism/index}{plagiarism} and thus an irregularity.
    \item Always check the translated text for correctness and meaning.
\end{itemize} \\ \hline
\textbf{As a search engine to get information on a topic or to search for existing research on the topic} & \textit{add text here} & This use is similar to e.g. a Google search or checking Wikipedia. If you write your own text based on this information, you do not have to refer to the use of GenAI in the text. \textbf{You only have to \href{https://bib.kuleuven.be/english/training-and-tutorials/citation}{refer} to the existing research and references you have checked and used }(without such references you would be committing \href{https://www.kuleuven.be/english/education/plagiarism/index}{plagiarism} and thus committing an irregularity).

\textbf{However, be careful:}
\begin{itemize}
    \item Be aware that the output of the GenAI tool cannot be guaranteed as a 100\% reliable source of information. The output may not be entirely correct and/or be limited due to the databases it uses. Moreover, knowledge evolves and may change over time; therefore, the database of the GenAI tool may not be up to date. Therefore, \textbf{verify the information and do not just copy-paste it} as you should understand and critically process everything you are writing.
\end{itemize} \\ \hline
\textbf{For literature search} & \textit{add text here} & This use is comparable to e.g.\ a Google Scholar search. You do not have to refer to such kind of GenAI use; \textbf{you only have to \href{https://bib.kuleuven.be/english/training-and-tutorials/citation}{refer} to the literature references you have checked and used} (without such references you would be committing \href{https://www.kuleuven.be/english/education/plagiarism/index}{plagiarism} and thus committing an irregularity).

\medskip
\textbf{However, be careful:}
\begin{itemize}
    \item Be aware that the search output is restricted to the database the GenAI tool is built on. After this initial search, look for scientific sources and conduct your own analysis of the source documents. \textbf{Interpret, analyse and process the information you obtained; verify it and do not just copy-paste it} as you should understand and critically process everything you are writing.
    \item Be aware that some GenAI tools may \textbf{output \href{https://www.kuleuven.be/english/education/leuvenlearninglab/support/toolguide/guidelines-for-safe-use-of-genai-tools}{no or wrong references.}} As a student you are responsible for further checking and verifying the absence or correctness of references; do not just copy-paste it.
\end{itemize} \\ \hline
\textbf{For generating programming code } & \textit{add text here} & Use of GenAI for coding is allowed except when it is prohibited by the teacher or the program of study. If used for coding, correctly mention the use of GenAI assistance in accordance with the instructions on the \href{https://www.kuleuven.be/english/education/student/educational-tools/generative-artificial-intelligence#Transparency}{page on being transparant about the use of GenAI.} \\ \hline
\textbf{For generating new (research) ideas} & \textit{add text here} & This use of GenAI is allowed except when it is prohibited by the teacher or the program of study. Correctly mention the use of GenAI assistance in accordance with the instructions on the \href{https://www.kuleuven.be/english/education/student/educational-tools/generative-artificial-intelligence#Transparency}{page on being transparent about the use of GenAI.}

\medskip
\textbf{Be careful:}
\begin{itemize}
    \item Further verify in this case whether the idea is novel or not. It is likely that it is related to existing work. \textbf{If so, that existing work should be correctly \href{https://bib.kuleuven.be/english/training-and-tutorials/citation}{referenced} in the text} (without such reference you would be committing \href{https://www.kuleuven.be/english/education/plagiarism/index}{plagiarism} and thus committing an irregularity).
\end{itemize}\\ \hline
\textbf{For generating synthetic data} & \textit{add text here} & Use of GenAI for generating synthetic data is allowed, \textbf{provided that it is methodologically and ethically justifiable}, except when it is prohibited by the teacher or the program of study. Always correctly mention the use of GenAI assistance in accordance with the instructions on the \href{https://www.kuleuven.be/english/education/student/educational-tools/generative-artificial-intelligence#Transparency}{page on being transparent about the use of GenAI.}

\medskip
\textbf{Be careful:}
\begin{itemize}
    \item Always \textbf{carefully evaluate} the generated synthetic data for quality and possible bias since the output is highly dependent on the quality of the data on which the models are trained.
\end{itemize} \\ \hline
\textbf{For generating blocks of text
(other than the allowed use without referencing mentioned above)
} &  & According to Article 84 of the Regulations on Education and Examinations your text should allow to correctly and properly assess your own knowledge, understanding and skills. \textbf{Therefore, inserting blocks of text without quotes and a reference to GenAI assistance in your work is not allowed. }

\medskip
\textbf{Be careful:}
\begin{itemize}
    \item If it is really needed to insert a block of text from a GenAI tool, for instance because of the nature of your assignment, mention it as a citation by using quotes and correctly mention the use of GenAI assistance in accordance with the instructions \href{https://www.kuleuven.be/english/education/student/educational-tools/generative-artificial-intelligence#Transparency}{page on being transparent about the use of GenAI.}
    \item However, in general, such GenAI use should be kept to an \textbf{absolute minimum}; you should always check the original sources.
\end{itemize} \\ \hline
\textbf{For generating visuals, video or audio} & \textit{add text here} & Use of GenAI for generating visuals, audio or video is allowed except when it is prohibited by the teacher or the program of study. 

\medskip
\textbf{Be careful:}
\begin{itemize}
    \item If used, refer to GenAI for visuals according to the usual referencing style, following the instructions on the \href{https://bib.kuleuven.be/english/training-and-tutorials/citation/refering-to-genai}{page on referencing GenAI.}
    \item If you work with existing visuals, audio or video, \textbf{that existing work should be correctly \href{https://bib.kuleuven.be/english/training-and-tutorials/citation}{referenced} in the text} (without such reference you would be committing \href{https://www.kuleuven.be/english/education/plagiarism/index}{plagiarism} and thus an irregularity).
    \item Explain the usage in the methods section (if there is one) and optionally attach (or link to) the prompts with the full output (or history).
\end{itemize} \\ \hline
\textbf{Other use (specify here; this may also include a combination of types of use mentioned above):} & \textit{add text here} & To make sure other use of GenAI is allowed within the course or thesis, and if so, the conditions that may apply, contact the teaching staff of the course or the supervisor of the thesis \textbf{beforehand} and explain the intended GenAI use. Also inform the programme director. 

\medskip
Motivate how you would comply with Article 84 of the exam regulations. Explain the use and the added value of the AI tool you consider to use and how it is in accordance with the assignment or thesis and the \href{https://www.kuleuven.be/english/genai}{KU Leuven guidelines on responsible use of Generative AI tools.} 

\medskip
Depending on the kind of GenAI use, it may be needed to properly reference it in the text, in accordance with the instructions \href{https://www.kuleuven.be/english/education/student/educational-tools/generative-artificial-intelligence#Transparency}{page on being transparent about the use of GenAI.}\\ \hline
\end{longtable}
\end{document}
